\subsubsection{Ternary congruence of the trace and optimality of the exponent}
\label{mg09:congruencia}

The level-three identity permits comparison of every coefficient of
the constructed trace with the corresponding coefficient of its modular
realization. The appropriate domain for this consequence is the ring
of formal power series: each coefficient depends on finitely many
factors of the product, while the statement retains all degrees.

\begin{lemma}[Integral unit of the order-three product]
The series
\begin{equation}
 u(q)=q\,t_3(q)=\prod_{n\geq1}(1+q^n+q^{2n})^{-12}
 \label{mg09:unidad}
\end{equation}
belongs to \(1+q\mathbb Z[[q]]\), and its inverse belongs to the
same set.
\end{lemma}

\begin{proof}
Each polynomial \(1+q^n+q^{2n}\) has constant term one. For any
series \(v(q)=1+\sum_{r\geq1}v_rq^r\) with integer
coefficients, the equation \(v(q)w(q)=1\) recursively determines
\[
 w_0=1,\qquad w_r=-\sum_{i=1}^{r}v_iw_{r-i}\quad(r\geq1).
\]
This recurrence proves existence, uniqueness and integrality of the
inverse. Applied to each factor, it also proves integrality of its
inverse twelfth power. Moreover,
\((1+q^n+q^{2n})^{-12}\in1+q^n\mathbb Z[[q]]\).
To compute a coefficient of degree at most \(r\), only factors
with \(n\leq r\) are needed; the partial products stabilize
coefficientwise. Consequently, \eqref{mg09:unidad} defines an
integral series with constant term one. Applying the same recurrence
to \(u\) yields \(u^{-1}\in1+q\mathbb Z[[q]]\).
\end{proof}

\begin{proposition}[Exact uniform divisibility]
For the trace \(t_3\) and the series \(J=j-744\) related by
\eqref{km:j-nivel3},
\begin{equation}
 J-(t_3+12)\in3^9q\mathbb Z[[q]].
 \label{mg09:congruencia-exacta}
\end{equation}
The uniform exponent of divisibility by three is exactly nine:
\begin{equation}
 [q]\bigl(J-t_3-12\bigr)=10\cdot3^9=196830.
 \label{mg09:optimalidad}
\end{equation}
\end{proposition}

\begin{proof}
Since \(t_3=q^{-1}u\), the lemma gives
\(t_3^{-r}=q^ru^{-r}\in q^r\mathbb Z[[q]]\) for
\(r=1,2,3\). Identity \eqref{km:j-nivel3} becomes
\begin{equation}
 J-t_3-12
 =3^9q\left(10u^{-1}+4\cdot3^5q\,u^{-2}
                         +3^9q^2u^{-3}\right).
 \label{mg09:factorizacion}
\end{equation}
Every term in parentheses is an integral series, proving
\eqref{mg09:congruencia-exacta}. Its constant term is \(10\),
since \(u(0)=u^{-1}(0)=1\), and the last two terms contain
\(q\) and \(q^2\). This gives \eqref{mg09:optimalidad}. As
\(3\nmid10\), the coefficient of \(q\) has \(3\)-adic valuation
nine, establishing optimality of the uniform exponent.
\end{proof}

In particular, \([q^r]J\equiv[q^r]t_3\pmod{3^9}\) for every
\(r\geq1\). The quantifier follows from the formal identity and
the integral inversion proved in the lemma. Verification through any
finite degree is a specialization of this relation. The trace continues
to arise from the order-three automorphism of the lattice constructed
from HMT incidence; uniformization subsequently acts on that trace.
The congruence preserves this order and determines an exact relation
between its two graded readings.
